
\subsection{MODELO CLIENTE-SERVIDOR}

A escolha dos componentes de um \en{software}, as interações entre eles e onde atuam, constitui a arquitetura de sistema. Na arquitetura de sistemas categoriza-se, basicamente, os sistemas em centralizados e descentralizados, essas categorizações se baseiam na natureza da relação cliente-servidor; Quando discutindo arquitetura de sistemas o modelo cliente-servidor é essencial. Desse modo, esse modelo define clientes que requisitam serviços de servidores, onde em sistemas centralizados é definido por: uma máquina cliente que implementa a parte do sistema voltada para o usuário; servidores contém o resto da aplicação (processamento e a camada de dados) \cite{tanenbaum2007dist}.

Em sistemas descentralizados essa organização difere, onde em sistemas centralizados a distribuição é alcançada, principalmente, por separar componentes logicamente em diferentes máquinas (distribuição vertical), nos descentralizados a distribuição é alcançada horizontalmente, isso é, clientes ou servidores podem ser divididos em partes equivalentes, cada uma responsável por processar sua parte do trabalho. Uma classe de arquitetura de sistema que implementa distribuição horizontal é \en{peer-to-peer} \cite{tanenbaum2007dist}.

\subsubsection{REDES \en{PEER-TO-PEER}}

Redes \en{peer-to-peer} são redes onde cada \en{node} atua como cliente e servidor. Essas redes podem ser de arquitetura estruturada e não estruturada. Nas estruturadas a rede é construída usando um procedimento determinista; um exemplo seria manter uma tabela com identificadores para todos os \en{nodes}, este mesmo identificador atua para categorizar os dados, assim quando um item é acessado na rede, é retornado o identificador do \en{node} responsável para ele. Em uma rede não estruturada os \en{nodes} conhecem uns aos outros de maneira quase aleatória, isso é, cada um mantém uma lista de vizinhos construída de maneira não determinista; a transmissão de dados ocorre da mesma forma \cite{tanenbaum2007dist}.

